\section{实验过程记录}

\subsection{问题1：ALU控制码与指导书不一致}
\textbf{问题描述：}发布包中 alu.v 的运算控制码与指导书要求不符，例如原代码中加法对应 op=010，而指导书要求 op=000，且原代码缺少 SLT 功能。

\textbf{解决方案：}按照指导书表1重新编写 alu.v，将6种运算的控制码对应关系修正，并使用 \texttt{\$signed()} 实现有符号比较完成 SLT 功能。zero 标志改为直接检测 result 是否为全0。

\subsection{问题2：仿真顶层设置错误导致波形全为X}
\textbf{问题描述：}在 Vivado 中运行仿真后，波形窗口所有信号均显示为 X 或 Z，无有效波形输出。

\textbf{解决方案：}检查发现仿真顶层被设置为 top.v（设计顶层）而非仿真文件 alu\_sim.v，top.v 中无激励信号导致输出未定义。将仿真顶层切换为 alu\_sim 后重新运行仿真，波形正常显示。

\subsection{问题3：流水线仿真文件使用时钟沿等待导致死循环}
\textbf{问题描述：}最初编写的流水线仿真文件在 for 循环内使用 \texttt{@(posedge clk)} 等待时钟沿，导致仿真时间异常（运行至2亿ns以上），仿真无法结束。

\textbf{解决方案：}ALU 为纯组合逻辑，无需等待时钟沿，将仿真激励改为使用 \texttt{\#10} 固定延时；流水线仿真文件中对 for 循环的时序控制进行重构，确保时钟与激励信号正确配合，仿真可正常结束。
